% =====================================================================
% Figure: TCDA Taxonomy (Section 3.1)
% This is a working draft. The user has indicated that final figure
% production will be handled separately; the file is provided so the
% document compiles and the cross-reference \Cref{fig:taxonomy} resolves.
% =====================================================================
\begin{figure}[t]
\centering
\begin{tikzpicture}[
  font=\small,
  level distance=11mm,
  level 1/.style={sibling distance=42mm, level distance=12mm,
                  every node/.style={draw, rounded corners, fill=blue!8,
                                     align=center, minimum width=32mm,
                                     minimum height=7mm, font=\small\bfseries}},
  level 2/.style={sibling distance=14mm, level distance=13mm,
                  every node/.style={draw, rounded corners, fill=teal!6,
                                     align=center, minimum width=28mm,
                                     minimum height=6mm, font=\footnotesize}},
  edge from parent/.style={draw, -{Latex[length=2mm]}, semithick}]

\node[draw, rounded corners, fill=violet!10, minimum width=64mm,
      minimum height=9mm, font=\bfseries] {Topological Causal Data Analysis (TCDA)}
  child {node {A.\,Topological\\Causal Representation}
    child {node {A1\\PH for Causal\\Structure Learning}}
    child {node {A2\\Mapper \& Reeb\\for Abstraction}}
    child {node {A3\\Causal Embeddings\\in Persistence}}
  }
  child {node {B.\,Topological\\Robustness}
    child {node {B1\\Stability of\\Causal Estimates}}
    child {node {B2\\Topological\\Confounding}}
    child {node {B3\\Functorial\\Causal Inference}}
  }
  child {node {C.\,Causal Interpretation\\of Topological Features}
    child {node {C1\\Betti Numbers as\\Causal Explanation}}
    child {node {C2\\Persistent Cohomology\\for Mediation}}
    child {node {C3\\Stratified Causal\\Heterogeneity}}
  }
  child {node {D.\,Applications\\\& Methodology}
    child {node {D1\\Time-Series\\Causality}}
    child {node {D2\\Spatial \&\\Spatiotemporal}}
    child {node {D3\\Industry\\Applications}}
  };
\end{tikzpicture}
\caption{The proposed TCDA taxonomy. Four orthogonal domains, each with
three sub-domains, organise the field. Domain A concerns
\emph{representation}: how topology summarises the structure of causal
data. Domain B concerns \emph{robustness}: how topological stability
controls the sensitivity of causal estimates to perturbations. Domain C
concerns \emph{interpretation}: which causal stories are encoded by
topological features. Domain D concerns \emph{methodology and applications}
in concrete data modalities. The taxonomy is intended to be exhaustive
(every existing TCDA contribution we are aware of fits in at least one
node) but not mutually exclusive (many works contribute to several
sub-domains simultaneously).}
\label{fig:taxonomy}
\end{figure}
